\section{半小映射的分解定理 (Xueqing Wen)}%
\label{sec:Decomposition theorem for semi-small maps (Xueqing Wen)}

\begin{quest}
    设 \(X\) 为（流形、簇、解析空间等）。我们如何计算 \(H^*(X, \C)\)？
\end{quest}

遗憾的是，这个问题太难了，因此我们将把它换成一个新的问题。

\begin{quest}
    我们如何利用更小的空间来探测 \(H^*(X, \C)\)？
\end{quest}

\begin{thm}[K\"unneth formula]
    我们有 \(H^*(X \times Y, \C) = H^*(X, \C) \otimes H^*(Y, \C)\)。
\end{thm}

遗憾的是，大多数空间都不是乘积，因此我们考虑如下问题。

\begin{quest}
    设 \(f \colon X \to Y\) 为态射。我们能否利用 \(Y\) 的上同调以及 \(f\) 的纤维的上同调来探测 \(X\) 的上同调？
\end{quest}

这个问题仍然非常困难，但至少我们可以研究它。我们将从 \(f\) 为紧流形之间的纤维丛的情形开始。答案当然是 Leray 谱序列或 Serre 谱序列，但在进入正题之前，我们需要一些预备知识。

\begin{defn}
    \(Y\) 上的 \(\ul{\C}_Y\)-局部系是指 \(Y\) 上向量空间的局部常值层 \(\mc{L}\)。
\end{defn}

例如，若 \(Y\) 连通，则局部系对应于 \(\pi_1(Y)\) 的表示。另请注意，若 \(Y\) 为代数簇，我们需要对其解析化进行研究。

\begin{exm}
    设 \(f\) 为如上以 \(F\) 为纤维的纤维丛。则 \(R^i f_* \ul{\C}_X\) 为 \(Y\) 上以 \(H^i(F, \C)\) 为纤维的局部系。
\end{exm}

\begin{prop}
    我们有
    \[ H^*(X, \C) = H^*(X, \ul{\C}_X) = H^*(\mr{pt}, R a_{X*} \ul{\C}_X). \]
\end{prop}

\begin{thm}[Leray/Serre spectral sequence]
    我们有谱序列
    \[ E_2^{p,q} = H^p(Y, R^q f_* \ul{\C}_X) \Rightarrow H^{p+q}(X, \C). \]
\end{thm}

我们现在考虑 \(f \colon X \to Y\) 为光滑射影簇之间的光滑射影映射的情形。

\begin{thm}[Deligne {\cite{Deligne1968}}]
    我们有非典范同构
    \[ R f_* \ul{\C}_X = \bigoplus_{q \geq 0} R^q f_* \ul{\C}[-q]. \]
    在上同调层面，我们有
    \[ H^*(X, \C) = \bigoplus_{q \geq 0} H^{*-q}(Y, R^q f_* \ul{\C}). \]
\end{thm}

这是射影簇的硬 Lefschetz 定理的推论，也表明 Serre 谱序列在 \(E_2\) 页退化。在拓扑中，这对 Hopf 纤维化不成立。

我们现在考虑 \(f \colon X \to Y\) 为代数簇之间的真映射的情形。此时存在 \textit{Whitney 分层}\footnote{我们明确拒绝定义它。} \(Y = \bigsqcup_{\lambda} Y_{\lambda}\)，使得 \(f^{-1}(Y_{\lambda}) \to Y_{\lambda}\) 为（可能奇异的簇的）\(C^{\infty}\) 纤维化。

\begin{defn}
    若 \(Y\) 上的 \(\ul{\C}_Y\)-模 \(\mc{F}\) 存在 Whitney 分层 \(Y = \bigsqcup Y_{\lambda}\)，使得对一切 \(\lambda\) 都有 \(\mc{F}|_{Y_{\lambda}}\) 为局部系，则称 \(\mc{F}\) 为\textit{可构造层}。
\end{defn}

有了这个定义，我们可以考虑 \(Y\) 上可构造层的导出范畴 \(\mc{D}_c^b(Y)\)。该构造具有六函子形式体系，其中函子为 \(f_*, f^*, f_!, f^!, \otimes, \mc{H}\mr{om}\)，并经由函子
\[ D_Y \colon \mc{D}_c^b(Y) \to \mc{D}_c^b(Y). \]
满足 Verdier 对偶。

\begin{defn}
    复形 \(\mc{F}^{\bullet} \in \mc{D}_c^b(Y)\) 若满足如下条件，则称为\textit{反常层}：
    \begin{enumerate}
        \item （支集条件）对一切 \(i\) 有 \(\dim \on{supp} \mc{H}^{-i}(\mc{F}^{\bullet}) \leq i\)；
        \item （余支集条件）对一切 \(i\) 有 \(\dim \on{supp} \mc{H}^{-i}(D_Y \mc{F}^{\bullet}) \leq i\)。
    \end{enumerate}
\end{defn}

\begin{rmks}\leavevmode
    \begin{enumerate}
        \item 反常层不是层，而是位于次数 \([-\dim Y, 0]\) 的复形。
        \item 它们有两个来源。其一是相交同调，其二是 Riemann-Hilbert 对应。
    \end{enumerate}
\end{rmks}

我们需要用到如下事实：\(Y\) 上反常层的范畴 \(\ms{Perv}(Y)\) 为 Abel 范畴，而且是 \(\mc{D}_c^b(Y)\) 上某个 \(t\)-结构的心。特别地，对任意 \(\mc{G} \in \mc{D}_c^b (Y)\)，我们可以考虑反常上同调 \(^p\mc{H}^i(\mc{G}) \in \ms{Perv}(Y)\)。

\begin{thm}[Beilinson-Bernstein-Deligne-Gabber {\cite{BBD1982}}]
    设 \(f \colon X \to Y\) 为代数簇之间的真态射。假设 \(X\) 光滑连通。\footnote{这不是必要的，但若去掉光滑性假设，我们需要考虑相交复形。}则
    \[ R f_* \ul{\C}_X [\dim X] = \bigoplus_{i} {^p\mc{H}^i}(R f_* \ul{\C}_X)[-i]. \]
    而且，每个 \(^p\mc{H}^i (R f_* \ul{\C}_X)\) 都是半单的（即每个子对象都是直和项）。
\end{thm}

\begin{defn}
    设 \(f \colon X \to Y\) 为真满态射，其中 \(X\) 光滑不可约。设 \(Y = \bigsqcup Y_{\lambda}\) 为分层，使得 \(f^{-1}(Y_{\lambda}) \to Y_{\lambda}\) 为 \(C^{\infty}\)-纤维化。若对一切 \(y_{\lambda} \in Y_{\lambda}\) 都有 \(\on{codim}_Y Y_{\lambda} \geq 2 \dim f^{-1} (y_{\lambda})\)，则称 \(f\) 为\textit{半小}的。若取等号，则称 \(Y_{\lambda}\) 为\textit{相关的}。若唯一相关的层为稠密层，则称 \(f\) 为\textit{小的}。
\end{defn}

\begin{thm}[{\cite{DeCataldoMigliorini2002}}]
    若 \(f \colon X \to Y\) 半小，则 \(R f_* \ul{\C}_X [\dim X]\) 为 \(Y\) 上的反常层。
\end{thm}

\begin{proof}
    对 \(y_{\lambda} \in Y_{\lambda} \subseteq Y\)，我们看到
    \[ R f_* \ul{\C}_X |_{y_{\lambda}} = \bigoplus H^i (f^{-1}(y_{\lambda}), \C). \]
    注意，若 \(i > 2 \dim f^{-1}(y_{\lambda}) \leq \on{codim}_Y (Y_{\lambda})\)，则此式为零。若让 \(Y_{\lambda}\) 取遍所有层，即得支集条件。因为 \(f\) 真，我们有 \(R f_* = f_!\)，再利用 \(D_Y R f_* = f_! D_X\) 即得余支集条件。
\end{proof}

\begin{exms}
    我们现在给出一些半小映射的例子。
    \begin{enumerate}
        \item 曲面之间的任何一般有限的真映射都是半小的。
        \item 设 \(Y = \ab\{xy = t^{\ell+1}\} \subseteq \C^3\)。则存在奇点的 crepant 消解 \(f \colon \tilde{Y} \to Y\)，而奇点上的纤维为一串 \(\P^1\)，故为半小的（亦见前一例子）。
    \end{enumerate}
\end{exms}

\begin{thm}[Kaledin {\cite{Kaledin2006}}]
    从非奇异全纯辛簇出发的射影双有理映射是半小的。
\end{thm}

\begin{exm}
    例如，Springer 消解和 Grothendieck-Springer 消解都是半小的。此外，曲面上点的 Hilbert 概形的 Hilbert-Chow 态射也是半小的。
\end{exm}

设 \(f \colon X \to Y\) 半小，并考虑分层 \(Y = \bigsqcup Y_{\lambda}\)。若 \(Y_{\lambda}\) 相关，考虑
\[ \mc{L}_{\lambda} = R^{d_{\lambda}} f_* \ul{\C}_{f^{-1}(Y_{\lambda})}, \]
其中 \(d_{\lambda} = 2 \dim f^{-1} (y_{\lambda})\)。则 \(\mc{L}_{\lambda}\) 半单。这里注意 \(\mc{L}_{|\lambda} = H^{d_{\lambda}}(f^{-1}(y_{\lambda}), \C)\)。现设 \(\mr{IC}(\mc{L}_{\lambda})\) 为 \(\ol{Y}_{\lambda}\) 上 \(\mc{L}_{\lambda}[\dim Y_{\lambda}]\) 的中间扩张，它是反常层。

\begin{thm}[{\cite{DeCataldoMigliorini2002}}]
    设 \(f \colon X \to Y = \bigsqcup Y_{\lambda}\) 半小。则存在典范分解
    \[ R f_* \ul{\C}_X[\dim X] = \bigoplus_{Y_{\lambda}\text{ relevant}} \mr{IC}(\mc{L}_{\lambda}). \]
\end{thm}

\begin{exm}
    设 \(G\) 为具 Lie 代数 \(\g\) 的约化群，并取 Borel 子群 \(B\)。则
    \[ G \times_B \b = \tilde{\g} \xrightarrow{\pi} \g \]
    是小的。因此我们有
    \[ R \pi_* \ul{\C}_{\tilde{\g}} [\dim \g] = \mr{IC}(\mc{L}). \]
    回忆 \(\tilde{\g}^{\on{rs}} \to \g^{\on{rs}}\) 为 \(W\)-torsor。则对 \(a \in \g^{\on{rs}}\)，我们有 \(\mc{L}_{|a} \cong \C[W]\)，从而 \(\mc{L} = R^0 \pi_* \ul{\C}_{\tilde{\g}^{\on{rs}}}\)。这蕴含 \(W\) 作用在 \(\mc{L}\) 上，此作用可延拓到 \(\mr{IC}(\mc{L})\) 上，也作用在 \(R \pi_* \C_{\tilde{\g}}\) 上，特别地作用在任意 \(a \in \g\) 的纤维 \(R \pi_* \C_{\tilde{\g}}|_a\) 上。
\end{exm}

\begin{exm}
    考虑 Springer 消解 \(\mu \colon T^*(\on{GL}_3/B) \to \mc{N} = \mathbb{O}_{(1^3)} \sqcup \mathbb{O}_{(2,1)} \sqcup \mathbb{O}_{(3)}\)。此映射半小，但遗憾的是所有幂零轨道都相关。注意 \((2,1)\) 轨道上的纤维为两条 \(\P^1\) 在节点处粘合而成，\((3)\) 轨道上的纤维为一点，而 \((1^3)\) 轨道上的纤维为 \(\on{GL}_3/B\)。

    我们看到 \(\mc{L}_{(3)}\) 必为平凡的，因为 \(\mu\) 在此轨道上为同构。接着 \(\mc{L}_{(2,1)} = R^2 \mu_* \ul{\C}\) 为 \(\mathbb{O}_{(2,1)}\) 上秩 \(2\) 的局部系。最后，\(\mc{L}_{(1^3)} = R^6 \mu_* \ul{\C}\) 为 \(\mathbb{O}_{(1^3)} = \mr{pt}\) 上秩 \(1\) 的局部系。因此我们有
    \[ R \mu_* \ul{\C}_{T^*\on{GL}_3/B}[6] = \mr{IC}(\mc{L}_{(3)}) \oplus \mr{IC}(\mc{L}_{(2,1)}) \oplus \mr{IC}(\mc{L}_{(1^3)}). \]
    限制到 \(b \in \mathbb{O}_{(3)}\)，得到 \(\mc{L}_{(3)} = H^0(\mr{pt})\)。限制到 \(a \in \mathbb{O}_{(2,1)}\)，得到
    \[ H^0(X, \C) \oplus H^2(X, \C) = \mr{IC}(\mc{L}_{(3)})|_{(a)} \oplus \mr{IC}(\mc{L}_{(2,1)})|_{(a)}. \]
    限制到 \(0 = \mathbb{O}_{(1^3)}\)，得到
    \begin{align*}
        H^*(\on{GL}_3/B, \C) = \C[W] &= \mr{IC}(\mc{L}_{(3)})|_0 \oplus \mr{IC}(\mc{L}_{(2,1)})|_0 \oplus \mc{L}_{(1^3)}|_0 \\
        &= H^0(\on{GL}_3/B, \C) \oplus \ab(\text{a summand of dimension } 4) \oplus H^6(\on{GL}_3/B, \C).
    \end{align*}
    作为 \(W = S_3\) 的表示，它们分别对应平凡表示、两份标准表示和符号表示。关于 Springer 理论的构造与这些分解，见~\cite{BorhoMacPherson1983}。
\end{exm}

\section{量子 K-理论 (Xiaohan Yan and You-Cheng Chou)}%
\label{sec:Quantum K-theory (Xiaohan Yan and You-Cheng Chou)}

\subsection{概述}%
\label{sub:Overview}

Gromov-Witten 不变量是“计数”光滑射影簇 \(X\) 中曲线的一种方式。GW 理论给出形如
\[ \ab<\alpha_1, \ldots, \alpha_n>_{g, n, \beta} \in \Q \]
的数，其中 \(\alpha_i \in H^*(X)\) 为各类，\(g\) 为亏格，\(n\) 为标记点个数，\(\beta \in H_2(X, \Z)\) 为曲线类。这已用于研究镜像对称，除镜像对称外还给出了双有理不变量。我们现在的目标是将整个理论提升到 K 理论。量子 K 理论给出数
\[ \ab<V_1, \ldots, V_n>_{g, n, \beta} \in \Z \]
其中 \(V_1, \ldots, V_n \in K(X)\)。它们满足各种差分方程\footnote{与上同调中的微分方程作比较。}，并在三维镜像对称与表示论中有多种应用。

\subsection{K-理论}%
\label{sub:K-theory}

\begin{defn}
    设 \(X\) 为光滑射影簇。则 \(X\) 的 \textit{K 理论} \(K(X)\) 为 \(X\) 上向量丛的 Grothendieck 群，关系为若
    \[ 0 \to V_1 \to V_2 \to V_3 \to 0 \]
    为短正合列，则 \(V_2 = V_1 + V_3\)。
\end{defn}

\begin{rmk}\leavevmode
    \begin{enumerate}
        \item \(K(X)\) 为环，运算为直和与张量积；
        \item Chern 特征给出同构 \(K(X) \to A^*(X)\)，至少在 \(\Q\) 上成立。
    \end{enumerate}
\end{rmk}

\begin{exm}
    设 \(X = \P^n\) 并记 \(P = \mc{O}(-1)\)。则我们有
    \[ K(X) = \C[P^{\pm 1}]/(P-1)^{n+1} \xrightarrow{\on{ch}} A^*(X) = \C[H]/H^{n+1}. \]
\end{exm}

对不假设光滑的 Deligne-Mumford 叠，有两个版本
\[ K_0(X) \supset K^0(X), \]
其中 \(K_0\) 指凝聚层，\(K^0\) 指完美复形。在光滑射影簇上二者重合，故从今往后 \(K(X)\) 指 \(K_0(X)\)。存在张量积
\[ K^0(X) \otimes K_0(X) \to K_0(X) \]
以及沿真态射的前推
\[ f_* \colon K(X) \to K(Y) \qquad V \mapsto \sum (-1)^i R^i f_* V \]
。若 \(f\) 为完美态射（例如平坦、正則嵌入等），则存在拉回
\[ f^* \colon K(Y) \to K(X) \qquad W \mapsto \sum (-1)^i \on{Tor}_i^{f^{-1} \mc{O}_Y} (\mc{O}_X, f^{-1} W). \]

\subsection{曲线的模空间}%
\label{sub:Moduli of curves}

回忆到 \(X\) 的稳定映射的模空间
\[ \bar{M}_{g,n}(X, \beta) = \ab\{ f \colon (\Sigma, p_1, \ldots, p_n) \to X \mid g(\Sigma) = g, f_*[\Sigma] = \beta, \abs{\Aut f} < \infty \} \]
。此空间具有赋值映射
\[ \ev_i \colon \bar{M}_{g,n}(X, \beta) \to X \qquad (f, \Sigma, \ul{p}) \mapsto f(p_i). \]
还有万有余切线 \(L_i = \sigma_i^* \omega_{ \mc{C}/\bar{M} }\)。

该模空间为真 DM 叠，具有完美阻碍理论，从而给出虚拟结构层 \(\mc{O}^{\vir}_{g,n,\beta}\)。它来自\textit{虚拟切丛}的结构，其中
\[ \mathbb{T}_{(f, \Sigma, \ul{p})} \bar{M}_{g,n}(X, \beta) = \chi(\Sigma, f^* T_X) - \chi(\Sigma, T_{\Sigma}(-p_1 - \cdots - p_n)). \]
粗略地说，若 \(\bar{M}\) 为簇 \(Y\) 上向量丛 \(E\) 的截面的零点轨迹，则 \(\mc{O}^{\vir} = \mc{O}_Y \otimes \mr{Eu}(E)\)，其中 K 理论 Euler 类为 \(\mr{Eu}(E) = \sum (-1)^i \Lambda^i E^{\vee}\)。

\subsection{KGW 不变量}%
\label{sub:KGW invariants}

关于量子 K 理论，见~\cite{Lee2004}；关于置换等变理论，见~\cite{Givental2017K1,GiventalK7}。To\"en~\cite{Toen1999} 证明了下面所用的 Deligne--Mumford 叠的 Riemann--Roch 公式。

\begin{defn}
    K 理论 GW 不变量定义为
    \[ \ab<V_1, \ldots, V_n>_{g,n,\beta} \coloneqq \chi\ab(\bar{M}_{g,n}(X, \beta); \mc{O}^{\vir} \otimes \bigotimes_{i=1}^n \ev_i^* V_i). \]
\end{defn}

若 \(V_1, \ldots, V_n\) 为真正的向量丛，则该不变量为整数，但一般我们考虑系数在特征零域中的 \(K_0\)。我们也可定义引力子代
\[ \ab<V_1 L_1^{k_1}, \ldots, V_n L_n^{k_n}>_{g,n,\beta} \coloneqq \chi\ab(\bar{M}_{g,n}(X, \beta); \mc{O}^{\vir} \otimes \bigotimes_{i=1}^n ( \ev_i^* V_i )L_i^{k_i} ). \]
利用 \(S_n\) 在 \(\bar{M}_{g,n}(X, \beta)\) 上由置换标记点给出的作用，则对象 \(\bigotimes_{i=1}^n \ev_i^* V\) 与 \(\bigotimes_{i=1}^n L_i^k\) 均为 \(S_n\)-等变的。

\begin{defn}
    我们定义\textit{置换不变 K 理论 GW 不变量}为
    \begin{align*}
        \ab<V L_1^k, \ldots, V L_n^k>_{g,n,\beta}^{S_n} &\coloneqq \chi\ab([ \bar{M}_{g,n}(X, \beta)/S_n ]; \mc{O}^{\vir} \otimes \bigotimes_{i=1}^n ( \ev_i^* V ) L_i^k) \\
        &= \dim \chi_{S_n} \ab(\bar{M}_{g,n}(X, \beta); \mc{O}^{\vir} \otimes \bigotimes_{i=1}^n ( \ev_i^* V ) L_i^k )^{S_n}.
    \end{align*}
\end{defn}

给定 \(t(q) = \sum_{n \in \Z} t_n q^n \in K(X)[q,q^{-1}]\)，可定义
\[ \ab< t(L_1), \ldots, t(L_n)>^{S_n}_{g,n,\beta} \coloneqq \chi ([\bar{M}_{g,n}(X,\beta)/S_n], \cdots). \]
更一般地，若 \(H = S_{a_1} \times \cdots \times S_{a_m} \subset S_n\)，可定义
\[ \ab<t_1(L_1), \ldots, t_1(L_{a_1}), t_2(L_{a_1+1}), \ldots, t_2(L_{a_1+a_2}), \ldots>_{g,n,\beta}^H. \]

\begin{rmk}
    我们想讨论与普通 GW 不变量的一些关系。回忆 Hirzebruch-Riemann-Roch 断言
    \[ \chi(M; V) = \int_M \on{ch}(V) \on{td}(T_M). \]
    在 \(M\) 为光滑 DM 叠的情形，T\"oen 证明了公式
    \[ \chi(M; V) = \int_{IM} \wt{\on{ch}}(V) \wt{\on{td}}(T_M), \]
    其中 \(IM = M \times_{\Delta, M \times M, \Delta} M\) 为 \(M\) 的惯性叠，波浪线表示取轨形 Chern 特征与 Todd 类。
\end{rmk}

\begin{exm}
    若 \(M = B\mu_r\)，则 \(IM = \bigsqcup_{g \in \mu_r} B \mu_r\)。若 \(M = \P(1, r)\) 为加权 \(\P^1\)，则 \(IM = M \sqcup \bigsqcup_{g \neq 1 \in \mu_r} B\mu_r\)。
\end{exm}

\begin{exm}
    设 \(M = BG\)，其中 \(G\) 为有限 Abel 群，\(V\) 为 \(BG\) 上向量丛，即 \(G\) 的表示。则 \(\chi(M, V) = \dim V^G\)。另一方面，我们有
    \[ \int_{IM} \wt{\on{ch}}(V) \cdot \wt{\on{td}}(T_M) = \int_{IM} \on{tr}_g V = \sum_{g \in G} \frac{1}{\abs{G}} \on{tr}_g V = \dim V^G, \]
    故在此情形验证了 Riemann-Roch。
\end{exm}

\subsection{公理性质}%
\label{sub:Axiomatic properties}

\begin{lem}[{\cite{GiventalK7}}]
    我们有\textit{弦方程}
    \begin{align*}
        \ab< 1, x_1(L), \ldots, x_k(L), t, \ldots, t >_{g, 1+k+n, \beta}^{S_n} = \ab<x_1(L), \ldots, x_k(L), t^{\otimes n}>^{S_n}_{g,k+n,\beta} \\
        + \sum_i \ab< \cdots, \frac{x_i(L) - x_i(1)}{L-1}, \ldots, t^{\otimes n}>^{S_n}_{g,k+n,\beta}.
    \end{align*}
\end{lem}

\begin{proof}
    证明与 GW 理论相同，来自沿映射 \(\pi \colon\bar{M}_{g,n+1}(X,\beta) \to \bar{M}_{g,n}(X,\beta)\) 计算 \(\pi^* L_i\)。
\end{proof}

\begin{lem}[{\cite{GiventalK7}}]
    我们有\textit{膨胀方程}
    \begin{align*}
        \ab<L-1, x_1(L), \ldots, x_k(L), t^{\otimes n}>_{g,1+k+n,\beta}^{S_n} ={}& (k-2)\ab<x_1(L), \ldots, x_k(L), t^{\otimes n}>_{g,k+n,\beta}^{S_n} \\
        &+ \ab<x_1(L), \ldots, x_k(L), t, t^{\otimes n-1}>_{g,k+n,\beta}^{S_{n-1}}.
    \end{align*}
\end{lem}

\begin{proof}
    我们利用事实
    \[ \pi_* L = k + \on{Ind}_{S_{n-1}}^{S_n}(1) - 1, \]
    其中 \(\pi\) 遗忘第一个标记点，而 \(L = \omega_{\pi} (p_1 + \cdots + p_n)\)。这来自存在留数映射
    \[ H^0(L) \xrightarrow{\on{Res}} \bigoplus_{i=1}^{k+n} \C p_i \xrightarrow{\sum} \C. \]
    中间项为 \(k\) 个平凡表示与一个标准表示的直和。
\end{proof}

\begin{notn}
    固定 \(\tau, t \in K(X)\)。记
    \[ \ab<\!\ab<x_1(L_1), \ldots, x_k(L_k)>\!>_{g,k} \coloneqq \sum_{m,n,\beta} \frac{Q^{\beta}}{m!} \ab<x_1(L_1), \ldots, x_k(L_k), \tau^{\otimes m}, t^{\otimes n}>_{g,k+m+n,\beta}^{S_n}. \]
    设 \(\{\phi_{\alpha}\}\) 为 \(K(X)\) 的基，\(\{\phi^{\alpha}\}\) 为对偶基。再定义配对
    \[ g_{\alpha\beta} = \chi(X, \phi_{\alpha} \otimes \phi_{\beta}). \]
    这还不是正确的配对，因此考虑
    \[ G_{\alpha\beta} = g_{\alpha\beta} + \ab<\!\ab<\alpha,\beta>\!>_{0,2}. \]
    设 \(g^{\alpha\beta} = (g_{\alpha\beta})^{-1}\) 与 \(G^{\alpha\beta} = (G_{\alpha\beta})^{-1}\)。
\end{notn}

\begin{lem}[{\cite{Givental2000K}}]
    给定 \(A,B,C,D \in K(X)[q,q^{-1}]\)，我们有 \textit{WDVV 方程}
    \begin{align*}
        \sum_{\alpha,\beta} \ab<\!\ab< A(L), B(L), \phi_{\alpha}>\! >_{0,3}& G^{\alpha\beta}\ab<\!\ab< C(L), D(L), \phi_{\beta}>\! >_{0,3} \\
        &= \sum_{\alpha,\beta} \ab<\!\ab< A(L), C(L), \phi_{\alpha}>\! >_{0,3} G^{\alpha\beta}\ab<\!\ab< B(L), D(L), \phi_{\beta}>\! >_{0,3}.
    \end{align*}
\end{lem}

\begin{proof}
    拉回 \(\bar{M}_{0,4} = \P^1\) 上边界除子的等式。在 K 理论中，当 \(D = \bigcup C_i\) 时，需要用容斥原理将 \(\mc{O}_D\) 与 \(\mc{O}_{C_i}\) 联系起来，这解释了 \(G^{\alpha\beta}\) 的出现。
\end{proof}

\subsection{环路空间形式}%
\label{sub:Loop space formalism}

\begin{defn}
    \textit{回路空间} \((\mc{K}, \Omega)\) 由
    \[ \mc{K} = K(X)(q) \ps{Q}\]
    与辛形式
    \[ \Omega(f, g) = - \on{Res}_{q = 0, \infty} \ab<f(q^{-1}), g(q)>_X \frac{\d{q}}{q}. \]
    组成。
\end{defn}

\begin{rmk}
    存在极化 \(\mc{K} = \mc{K}_+ \oplus \mc{K}_-\)，其中
    \[ \mc{K}_+ = K(X)[q, q^{-1}]\ps{Q} \]
    而
    \[ \mc{K}_- = \ab\{ \frac{f}{g} \mid f, g \in K(X)\ps{Q}[q], g(0) \neq 0, \deg f < \deg g\}. \]
    例如，\(\frac{1}{1-q} \in \mc{K}_-\)。
\end{rmk}

\begin{defn}
    \textit{大 J 函数} \(J \colon \mc{K}_+ \to \mc{K}\) 由下式给出
    \[ \mathbf{t} = \mathbf{t}(q) \mapsto 1 - q + \mathbf{t}(q) + \sum_{\beta, \alpha} Q^{\beta} \phi^{\alpha} \ab< \frac{\phi_{\alpha}}{1-qL}, \mathbf{t}(L)^{\otimes n}>_{0, 1+n, \beta}^{S_n}. \]
\end{defn}

\begin{rmk}
    注意应用 Riemann-Roch，我们有
    \begin{align*}
        \chi\ab([\bar{M}/S_n], \frac{1}{1-qL}) &= \int_{I[\bar{M}/S_n]} \on{ch}\ab(\frac{1}{1-qL}) \on{td}.
    \end{align*}
    每一项都形如 \(\on{ch}\ab(\on{tr}_g \frac{1}{1-qL})\)，从而给出形如 \(\on{ch}\ab(\frac{1}{1-q\zeta L})\) 的项，其中 \(\zeta\) 为某单位根。再展开
    \[ \frac{1}{1-q\zeta L} = \frac{1}{(1-q \zeta) - q\zeta (L-1)} = \sum_k \frac{(q\zeta)^k}{(1-q\zeta)^{k+1}}(L-1)^k. \]
\end{rmk}

\begin{defn}
    \textit{Givental 锥}定义为 \(\mc{L} \coloneqq \Ima J \subset \mc{K}\)，其中我们在 \(Q\)-adic 拓扑下考虑 \(1-q\) 附近的芽。
\end{defn}

\begin{exm}
    设 \(X = \P^n\)。则看到
    \[ J(0) = (1-q) \sum_{d \geq 0} \frac{Q^d}{\prod_{\ell=1}^d (1-Pq^{\ell})^{n+1}}, \]
    其中 \(P = \mc{O}(-1)\)。如前可展开
    \[ \frac{1}{1-P q^{\ell}} = \frac{1}{(1-q^{\ell}) - q^{\ell}(P-1)} = \sum_{k=0}^{\infty} \frac{q^{\ell k}}{(1-q^{\ell})^{k+1}}(P-1)^k. \]
\end{exm}

\begin{prop}[{\cite{GiventalK7}}]
    事实上，\(\mc{L}\) 为 \(\mc{K}\) 中的锥，且我们有
    \[ \mc{L} = \bigcup_{t \in K(X)\ps{Q}} (1-q) S_{t}^{-1}(q) \mc{K}_+. \]
\end{prop}

我们也可引入“巨大”J 函数
\[ J(\mathbf{t}, \tau) = 1-q + \mathbf{t}(q) + \tau + \sum_{\beta, k, \alpha} \frac{Q^{\beta}}{k!} \phi^{\alpha} \ab< \frac{\phi_{\alpha}}{1-qL}, \mathbf{t}(L)^{\otimes n}, \tau^{\otimes k}>_{0,1+n+k,\beta}^{S_n}. \]
若固定 \(\mathbf{t}(q)\)，则 \(J(\tau)\) 来自 \(\tau\) 的 Frobenius 流形上\textit{微分方程}的基本解。此外，\(\mc{L}^{\mathbf{t}}=  \Ima J(\tau)\) 实际上为 Lagrangian 锥。

\begin{exm}
    考虑 \(f \colon \C[\Lambda_1^{\pm}, \Lambda_2^{\pm}] \to \C[\Lambda_1^{\pm}, \Lambda_2^{\pm}]\)，由 Adams 运算给出
    \[ t \mapsto \on{tr}_{(1\cdots n)} t^{\otimes n} = \Psi^n (t). \]
    因此它将 \(\Lambda_1 \mapsto \Lambda_1^n\) 与 \(\Lambda_2 \mapsto \Lambda_2^n\)。然而不清楚如何在 \(t\) 中求导，因为需要 \(nv t^{n-1} = \Psi^n(v)\)，故 J 函数在 \(\mathbf{t}\) 中不满足量子微分方程。
\end{exm}

\subsection{等变 K-理论与局部化}%
\label{sub:Equivariant K-theory and localization}

关于不动点递推与射影空间计算，见~\cite{GiventalK2}。

设环面 \(T = (\C^{\times})^r\) 作用在光滑射影簇 \(X\) 上。

\begin{defn}
    \textit{等变 K 理论} \(K_T(X)\) 为 \(X\) 上 \(T\)-等变向量丛的 Grothendieck 群。
\end{defn}

\begin{exm}
    例如我们有
    \[ K_T(\on{pt}) = \on{Rep}(T) = \C[\Lambda_1^{\pm}, \ldots, \Lambda_r^{\pm}], \]
    其中 \(\Lambda_i\) 为权 \((0, \ldots, 1, \ldots, 0)\) 的 \(1\) 维表示。
\end{exm}

\begin{rmk}
    对任意 \(X\)，\(K_T(X)\) 为 \(K_T(\mr{pt})\)-代数。
\end{rmk}

\begin{exm}
    考虑 \(\P^n\) 的完全等变 K 理论。则有
    \[ K_T(\P^n) = \C[P^{\pm}, \Lambda_0^{\pm}, \ldots, \Lambda_n^{\pm}] \big/ \prod_{i=0}^n (P-\Lambda_i). \]
    当 specialization \(\Lambda_i = 1\) 时，回到非等变 K 理论。
\end{exm}

\begin{defn}
    可定义等变 Euler 示性数为前推
    \[ \chi_T = \pi_* \colon K_T(X) \to K_T(\mr{pt}). \]
\end{defn}

\begin{exm}
    例如有 \(\chi_T(\mc{O}_{\P^1}(1)) = \Lambda_0^{-1} + \Lambda_1^{-1}\)。类似地有
    \[ \chi_T(\mc{O}(2)) = \Lambda_0^{-2} + \Lambda_0^{-1} \Lambda_1^{-1} + \Lambda_1^{-2}. \]
\end{exm}

\begin{notn}
    注意 \(\on{Frac}(T) = \on{Frac}(\on{Rep}(T))\)。再记
    \[ K_T^{\loc}(X) \coloneqq \on{Frac}(T) \otimes_{\on{Rep}(T)} K_T(X). \]
\end{notn}

我们需要如下事实。设 \(F\) 为 \(T\) 在 \(X\) 上的不动点集，并考虑 \(\iota \colon F \hookrightarrow X\)。则存在 \(\on{Frac}(T)\)-代数的同构
\[ K_T^{\loc}(X) \xrightarrow{\iota^*} K_T^{\loc}(F). \]

\begin{rmk}
    其逆由下式给出
    \[ (\iota^*)^{-1} = \bigoplus_{i=1}^r \frac{\iota_{i*}}{\on{Eu}(N_{F_i/X})}, \]
    其中 \(\on{Eu}(V) \coloneqq \sum_{j = 0}^{\on{rk} V} (-1)^j \Lambda^j V^*\) 为 K 理论 Euler 类。
\end{rmk}

\begin{exm}
    在 \(\P^n\) 上有 \(n+1\) 个孤立不动点 \(F_i\)，于是得到
    \[ \frac{\iota_{i*} \mc{O}_{F_i}}{\on{Eu}(N_{F_i/\P^n})} = \prod_{j \neq i} \frac{P - \Lambda_j}{\Lambda_i - \Lambda_j}. \]
    它们为 \(K_T^{\loc}(\P^n)\) 的幂等元。
\end{exm}

利用等变局部化得到
\[ \chi_T(X; V) = \sum_i \chi\ab(F_i; \frac{\iota_i^* V}{\on{Eu}(N_{F_i/X})}). \]

\begin{rmk}
    在 \(\Lambda = 1\) 极限下可回到非等变 Euler 示性数。
\end{rmk}

\begin{exm}
    设 \(X = \P^1\)。则计算得
    \begin{align*}
        \chi_T(\P^1, \mc{O}(1)) &= \chi\ab(F_0, \frac{\mc{O}(1)|_{F_0}}{\on{Eu}(T_{F_0/\P^1})}) + \chi\ab(F_1, \frac{\mc{O}(1)|_{F_1}}{\on{Eu}(T_{F_1/\P^1})}) \\
        &= \frac{\Lambda_0^{-1}}{1 - \frac{\Lambda_0}{\Lambda_1}} + \frac{\Lambda_1^{-1}}{1 - \frac{\Lambda_1}{\Lambda_0}} \\
        &= \Lambda_0^{-1} + \Lambda_1^{-1}.
    \end{align*}
\end{exm}

从今往后假设 \(X\) 为 GKM 流形。我们现在的目标是用 \(\bar{M}_{g,n}(X, \beta)\) 上的局部化来研究 \(X\) 的大 J 函数。任何 \(T\)-不动稳定映射 \(f \colon (\Sigma, \ul{p}) \to X\) 有两类不可约分支。一类是整个映到不动点的分支，另一类是经由一维轨道的分歧覆盖映射的分支。遗憾的是我们不能经由虚拟局部化计算大 J 函数，但仍可研究它。

设 \(f = f(q) \in \mc{K} = K_T^{\loc}(X)(q)\ps{Q}\)。记不动点为 \(F_{\alpha}\)（\(\alpha = 1, \ldots, r\)），包含映射为 \(\iota_{\alpha} \colon F_{\alpha} \hookrightarrow X\)。再记
\[ f_{\alpha}(q) = \iota_{\alpha}^* f(q) \in K_T^{\loc}(\on{pt})(q)\ps{Q}. \]
则 \(f \in \mc{L}\) 当且仅当对一切 \(\alpha\) 如下条件成立：
\begin{enumerate}
    \item 给定 \(\lambda \in T_{F_{\alpha}} X\) 与 \(m \in \N_+\)，\(f_{\alpha}(q)\) 至多在 \(q = \lambda^{1/m}\) 处有单极点，且
        \[ \on{Res}_{q = \lambda^{1/m}} f_{\alpha}(q) \frac{\d{q}}{q} = -\frac{Q^m}{m} \cdot \frac{\on{Eu}(T_{F_{\alpha}/X})}{\on{Eu}(T^{\vir}_{\phi_{\alpha,\alpha',m}}\bar{M}_{0,2}(X, m \cdot \beta(\alpha, \alpha')))} \cdot f_{\alpha'}(\lambda^{1/m}). \]
        这里，\(\lambda \in T_{F_{\alpha}}X\) 指 \(\lambda\) 为 \(T_{F_{\alpha}} X\) 的 \(1\) 维子表示。而且，\(\beta(\alpha,\alpha')\) 为由权 \(\lambda\) 决定的 \(T\) 的一维轨道的次数，\(\phi_{\alpha, \alpha', m}\) 参数化该一维轨道的 \(m:1\) 覆盖。
    \item \(f_{\alpha}(q)\) 的其余极点在 \(0, \infty\) 与单位根处，其余处正则。此外还有各种细小条件。
\end{enumerate}

\begin{exm}
    设 \(X = \P^n\)。则
    \[ J_T^X(0) = (1-q) \sum_{d \geq 0} \frac{Q^d}{\prod_{i=0}^n \prod_{\ell = 1}^d \ab(1 - \frac{P}{\Lambda_i} q^{\ell})} \in K_T(X)(q) \ps{Q}. \]
    于是蕴含
    \[ J_T^X(0) |_{F_j} = (1-q) \sum_{d \geq 0} \frac{Q^d}{\prod_{i=0}^n \prod_{\ell = 1}^d \ab(1 - \frac{\Lambda_j}{\Lambda_i}q^{\ell})}, \]
    故极点在 \(q = \ab(\frac{\Lambda_i}{\Lambda_j})^{1/\ell}\)。另注意
    \[ T_{F_j} X = \frac{\Lambda_0}{\Lambda_j} + \cdots + \frac{\Lambda_n}{\Lambda_j} - 1. \]
\end{exm}

\subsection{量子 K/GV 对应}%
\label{sub:Quantum K/GV correspondence}

quintic 对应猜测于~\cite{GaroufalidisScheidegger2022}，证明于~\cite{ChouLeeQuintic}。如下 Calabi--Yau 三维形公式属于 Chou--Lee~\cite{ChouLee2025}。

我们这里的目标是在 Calabi-Yau 三维形上计算 \(J^K(0)\)。如下结果由 Jockers-Mayr 与 Garoufalidis-Scheidegger 猜测。

\begin{thm}[{\cite{ChouLee2025}}]
    我们有
    \begin{align*}
        J^K(0) &\coloneqq (1-q) + \sum_{\beta, \alpha} Q^{\beta} \Phi^{\alpha} \ab<\frac{\Phi_{\alpha}}{1-qL}>_{0,1,\beta}^{X,K} \\
        &= (1-q) + \sum_{\beta} \sum_{r=1}^{\infty} \mr{GV}_{\beta} Q^{r\beta} \ab[\sum_{j=1}^{n_1} \Phi^{1j} a(r, q^r) + \Phi^{01} b(r, q^r)],
    \end{align*}
    其中 \(\ab\{\Phi_{\alpha}\}_{\alpha=1}^N\) 为 \(K^0(X)\) 的基，分解为
    \[ \ab\{\Phi_{\alpha}\}_{\alpha=1}^N = \bigsqcup_{i=0}^3 \ab\{\Phi_{ij}\}_{j=1}^{n_i} \]
    使得固定 \(i\) 时 \(\ch(\Phi_{ij})\) 的全体构成 \(H^{2i}(X)\) 的基。再取对偶基 \(\ab\{\Phi^{\alpha}\}_{\alpha} = \ab\{\Phi^{ij}\}_{i,j}\)。而且设多重覆盖公式为
    \begin{align*}
        a(r, q^r) &= \frac{r-1}{1-q^r} + \frac{1}{(1-q^r)^2}, \\
        b(r, q^r) &= \frac{r^2-1}{1-q^r} + \frac{3}{(1-q^r)^2} - \frac{2}{(1-q^r)^3}.
    \end{align*}
    最后，GV 即 Gopakumar-Vafa 不变量，在亏格零由下式定义
    \[ \mr{GW}_{0,\beta} = \sum_{r \mid \beta} \frac{1}{r^3} \mr{GV}_{\beta/r}. \]
    由 M\"obius 反演看到
    \[ \mr{GV}_{\beta} = \sum_{r \mid \beta} \frac{\mu(r)}{r^3} \mr{GW}_{0,\beta/r}. \]
\end{thm}

\begin{cor}[{\cite{ChouLee2025}}]
    我们有
    \[ \ab<1>_{0,1,\beta}^K = \sum_{r \mid \beta} r^2 \mr{GV}_{\beta/r}. \]
    应用 M\"obius 反演得到\footnote{注意此结果在一切亏格由 Ionel--Parker~\cite{IonelParker2018} 得到。}
    \[ \mr{GV}_{\beta} = \sum_{r \mid \beta} r^2 \mu(r) \ab<1>_{0,1,\beta/r}^K \in \Z. \]
    此外，可经由下式回到通常的 GW 不变量\footnote{这就像若应用 Hirzebruch-Riemann-Roch 计算 K 理论配对，最高次项只看到两个输入的 Chern 特征。}
    \[ \lim_{q \to 1} (1-q)^3 \ab<\frac{1}{1-qL}>_{0,1,\beta}^{X,K} = \sum_{r \mid \beta} \frac{-2}{r^3} \mr{GV}_{\beta/r} = -2 \mr{GW}_{\beta} = \ab<\psi>_{0,1,\beta}^{\mr{GW}}. \]
\end{cor}

粗略地可认为 \(J^K(0) \in \Z\ps{q}\ps{Q}\)，但更准确地说可考虑
\[ J^K(0) \in 1-q + \mc{K}_-. \]
这里展开 \(\frac{1}{1-qL} = \sum q^i L^i\)。
我们还需考虑如下事实：若 \(M\) 为 DM 叠而 \(L\) 为 \(M\) 上线丛，则函数
\[ f^{\Phi}(q) \coloneqq \chi\ab(M, \frac{\Phi}{1-qL}) \in \Z \ps{q} \]
在 \(\abs{q} < \ep\) 时收敛为有理函数。注意极点仅在单位根处。

此事实可用轨形 Riemann-Roch 定理证明。将惯性分解为
\begin{align*}
    IM &= \bigsqcup_{i \in I} M_i \\
    &= \bigsqcup_{\zeta} M_{\zeta},
\end{align*}
其中 \(M_{\zeta}\) 为扭 sector 的并，使得 \(g\) 在 \(L\) 上的作用的特征值为 \(\zeta\)。利用此可计算
\begin{align*}
    \chi\ab(M, \frac{1}{1-qL}) &= \sum_{i \in I} \int_{M_i} \on{td}(M_i) \on{ch} \ab(\frac{ \on{Tr}_g \ab(\frac{1}{1-qL}) }{\lambda_{-1}(N_{M_i/M}^*)}) \\
    &= \sum_{\zeta} \int_{M_{\zeta}} \on{td}(M_{\zeta}) \sum_{k \geq 0} \frac{(\zeta q)^k}{(1-\zeta q)^{k+1}} \frac{( \on{ch}(L)-1 )^k}{\ch \lambda_{-1}(N^*)},
\end{align*}
故一切极点必在单位根处。这里用展开式
\[ \frac{1}{1-\zeta q L} = \sum_{k \geq 0} \frac{(\zeta q)^k}{(1-\zeta q)^{k+1}} (L-1)^k \]
并取 Chern 特征。

\begin{exm}
    设 \(L = \mc{O}(1)\) 在 \(\P^1\) 上。则有
    \begin{align*}
        \chi\ab(\P^1, \frac{1}{1-qL}) &= \chi\ab(\P^1, \frac{1}{1-q} + \frac{q}{(1-q)^2}(L-1)) \\
        &= \sum_{i \geq 0} q^i (i+1) \\
        &= \frac{1}{(1-q)^2}.
    \end{align*}
\end{exm}

\begin{exm}
    我们现在考虑加权射影空间。则有
    \begin{align*}
        \chi\ab(\P(w_0, \ldots, w_n), \frac{1}{1-qL}) &= \prod_{i=0}^n \frac{1}{1-q^{w_i}}.
    \end{align*}
\end{exm}
